\chapter{Classical Theories}

Consider the linear beam theories developed in \cref{ch:frame-hierarchy}.

\section{Euler}

In order to diagonalize, one only needs to adopt a coordinate system originating at the \emph{shear center}, and oriented along the principal axes. This is because the only coupling that arises is due to the \(\LinearMW\) term.
%
This is not sufficient for the shear formulation. However, if torsional warping can be considered uniform (ie, \(\alpha' = 0\)) then the equations are nearly diagonalized by coordinates originating at the \emph{centroid}.

\begin{equation*}
\bm{e} = \begin{pmatrix} 
u_{x}^{\prime} \\ 
\vartheta^{\prime} \\ 
-u_z^{\prime\prime} \\ 
 u_y^{\prime\prime} \\
\end{pmatrix} 
\end{equation*}


\begin{equation*}
\begin{aligned}
N &= EA u_{x}^{\prime} \\ 
T &= G I \vartheta^{\prime} - G J_v \alpha \\ 
M_y &= -EI_y u_z^{\prime\prime} \\ 
M_z &=  EI_z u_y^{\prime\prime} \\
B   &= E J_w \alpha' \\
Q   &= G J_v ( \alpha - \vartheta' )\\
\end{aligned} 
\end{equation*}

$$
\beta = \sqrt{I_{y}^{2} - 2 I_{y} I_{z} + 4 I_{yz}^{2} + I_{z}^{2}}
= \sqrt{(I_y - I_z)^2 + 4 I_{yz}^2}
$$
$$
\frac{1}{2}(I_{0}-\beta),  
\begin{pmatrix}
\frac{1}{I_{yz}} \left(I_{z}- \frac{1}{2}( I_{0} - \beta)\right)\\
1
\end{pmatrix}
$$
$$
\frac{I_{0}}{2} + \frac{\beta}{2}, 
%
\begin{pmatrix}
\frac{1}{I_{yz}}
\left(
I_{z} -\frac{1}{2} \left(
  I_{0} + \beta
\right)
\right) \\
1
\end{pmatrix}
%
$$

\subsection{Type 1}

and equilibrium:
\begin{equation*}
\begin{aligned}
EA u_{x}'' &= 0 \\ 
(GI \vartheta' - G J_v  \alpha)'  &= 0 \\ 
-EI_y u_z'''  &= 0 \\ 
 EI_z u_y'''  &= 0 \\
E J_w \alpha'' - G J_v  ( \alpha - \vartheta' ) &= 0 
\end{aligned} 
\end{equation*}

\[
\renewcommand{\arraystretch}{1.2}
\begin{NiceArray}{ @{} l @{} >{{}}c<{{}} @{} l l @{} }
% \delta\psi   &=&  i\beta\gamma^{5} \\
EA u_{x}'' &=& 0 & \Block{2-1}{\operatorname{B}_{\mathrm{p}}} \\ 
(GI \vartheta' - G J_v  \alpha)'  &=& 0 \\ 
-EI_y u_z'''  &=& 0 \\ 
 EI_z u_y'''  &=& 0 \\
 E J_w \alpha'' - G J_v  ( \alpha - \vartheta' ) &=& 0 
\CodeAfter\SubMatrix.{1-1}{4-3}\}
\end{NiceArray}
\]
\subsection{Type 2}

Take \(\alpha = \vartheta'\) so that \(B = E J_w \vartheta''\). With regard to \cref{eq:type2-equilibrium}, the equilibrium equation for torsion is composed of the terms:
\begin{equation*}
\bigl(\underbrace{(GI - G J_v ) \vartheta'}_{\bar{\bm{s}}_{\mathrm{p}}} -  \underbrace{E J_w \vartheta'''}_{\mathbb{P}^*\bm{w}'}\bigr)'
\end{equation*}
thus furnishing the classical equations
\begin{equation*}
\begin{aligned} 
EA u_{x}'' &= 0 \\ 
\bigl((GI - G J_v ) \vartheta' -  E J_w \vartheta'''\bigr)'  &= 0 \\ 
-EI_y u_z''' &= 0 \\
 EI_z u_y''' &= 0 \\
\end{aligned} 
\end{equation*}

\subsection{Type 3}
In this case we further assume \(\alpha' = 0 \; (=\vartheta'')\)
\begin{equation*}
\begin{aligned} 
EA u_{x}'' &= 0 \\ 
\bigl(G(I - J_v ) \vartheta' \bigr)'  &= 0 \\ 
-EI_y u_z''' &= 0 \\
 EI_z u_y''' &= 0 \\
\end{aligned} 
\end{equation*}

\section{Shear}


\subsection{Type 1}


\subsection{Type 2}


\subsection{Type 3}

